\section{Where the backdoor lives in a ViT}
\label{sec:mechanism}

\Cref{sec:results} is a ranking. This section explains it by following a backdoor through a ViT, each experiment on the same backdoored models with the benign model of the same dataset, probed with the same trigger, as the control.

\subsection{A direction, written late}
\label{sec:mechanism:direction}

For a model and a layer, the backdoor direction is the mean difference between the class token's activation on a triggered image and on its clean copy, the construction Karayal\c{c}in et al. use on ViTs. \Cref{tab:direction-ablation} in \cref{app:mechanism} tests whether it is causal. Projecting the final layer's direction out of the class token takes attack success to \AblationDirectionMinAsr{} on BadNets and to at most \AblationDirectionMaxAsr{} on any of the \AblationCheckpoints{} attacks, while clean accuracy falls by at most \AblationCleanAccuracyCost{}. Projecting out a random direction of the same norm leaves attack success at \AblationRandomDirectionMinAsr{} or above, and so does zeroing the 20 coordinates the trigger moves most. The backdoor is a direction in the residual stream rather than a set of neurons, so a method that ranks coordinates, heads or channels sees only a rotated shadow of it. LF, the exception at the final layer, falls to \ErasureLfAsrAfterAllWritesMax{} once the direction is removed from every block's write to the stream. The direction stays near the benign level through the first half of the network and is written in the last third (\cref{app:mechanism}).


\subsection{Routed through attention from the trigger's own tokens}
\label{sec:mechanism:routing}

Activation patching \cite{zhang2024patching} asks where the backdoor sits at each depth. At a chosen block we overwrite a group of tokens on the triggered image with their clean values and measure how much of the clean prediction comes back (\cref{fig:patching} in \cref{app:mechanism}). For the local triggers, patching the trigger's own tokens recovers \PatchingLocalEarlyTrigger{} of the clean prediction over blocks \PatchingEarlyLayers{}, and the class token overtakes them only at block \PatchingLocalCrossover{} on average. A random group of the same size never exceeds \PatchingRandomNullMax{}. The class token's attention on the trigger's tokens rises from \RoutingBadnetClsEarly{} in blocks 1 to 4 to \RoutingBadnetClsLate{} in blocks 9 to 12 on BadNets, against \RoutingBenignClsLate{} on the benign model. For a local trigger the backdoor therefore sits in the trigger's own tokens for most of the network and reaches the class token through attention only in the last quarter. A global trigger covers every patch, so this test can only confirm that its backdoor is spread over the tokens.


\subsection{Which predictions survive each placement}
\label{sec:mechanism:placements}

PSBD flags an input whose prediction survives the perturbation, so a placement works when it changes clean predictions and leaves triggered ones in place. Every sweep records both shares, and we read them at each placement's adaptive rate, where PSBD-TM changes \SurvivalTmClean{} of clean test predictions and PSBD-RD \SurvivalRdClean{} (\cref{fig:survival}, per attack in \cref{tab:survival}). PSBD-TM changes only \SurvivalTmBadnetATwooTriggered{} of triggered BadNets predictions, \SurvivalTmBlendTriggered{} of Blend and \SurvivalTmLfTriggered{} of LF. PSBD-RD leaves the global triggers almost as intact, changing \SurvivalRdGlobalTriggered{} of their triggered predictions against \SurvivalTmGlobalTriggered{} for PSBD-TM, but it changes \SurvivalRdBadnetATwooTriggered{} of triggered BadNets predictions. The two placements therefore part on patch triggers. BadNets reads \SurvivalTmBadnetATwooAuroc{} under PSBD-TM and \SurvivalRdBadnetATwooAuroc{} under PSBD-RD, while the global triggers read \SurvivalTmGlobalAuroc{} and \SurvivalRdGlobalAuroc{}.

\begin{figure*}[t]
\centering
\includegraphics[width=\textwidth]{mech_survival.pdf}
\caption{Share of triggered predictions that change against the share of clean predictions that change, over the swept rates, averaged over the models of each attack. The vertical line is the adaptive target and the dotted diagonal is no separation. Patch triggers are solid and global triggers dashed. Under PSBD-TM the BadNets curve stays near zero until almost every clean prediction has changed, and under PSBD-RD it follows the diagonal.}
\label{fig:survival}
\end{figure*}

\Cref{tab:survival-sites} in \cref{app:survival} asks which property of a perturbation lets a patch trigger through, on the \SurvivalSiteModels{} BadNets models. Masking whole tokens at the attention input changes \SurvivalSiteMaskAttentionInputTriggered{} of triggered predictions and the same mask on the attention branch output changes \SurvivalSiteMaskAttentionOutputTriggered{}. Applied to the residual stream after the attention addition, the same operator changes \SurvivalSiteMaskStreamTriggered{}, close to the \SurvivalSiteDropoutStreamTriggered{} of PSBD-RD. At the attention input, dropout changes \SurvivalSiteDropoutAttentionInputTriggered{}, channel masking \SurvivalSiteChannelAttentionInputTriggered{} and Gaussian noise \SurvivalSiteNoiseAttentionInputTriggered{}. All three perturb every token instead of removing some tokens whole. The perturbations that leave a patch trigger's prediction in place are the ones that leave the residual stream untouched and remove whole tokens rather than perturbing every token. That fits where \cref{sec:mechanism:routing} places the backdoor, in the trigger's own tokens of the residual stream until the last blocks. A mask on a sublayer input or a branch output takes a token out of one block's computation and leaves its entry in the stream for the next block, while a perturbation of the stream changes that entry itself. These readings show which perturbations leave the trigger intact. They do not by themselves isolate which of the two properties does the work.


TaCT is the exception among patch triggers. Its trigger acts only on images of the source class, and PSBD-TM changes \SurvivalTmTactTriggered{} of its triggered predictions against \SurvivalTmClean{} of clean ones, and PSBD-RD changes \SurvivalRdTactTriggered{}. PSBD-TM still separates it at \SurvivalTmTactAuroc{} AUROC, and PSBD-RD reads \SurvivalRdTactAuroc{}.

\subsection{LayerNorm absorbs noise and not masking}
\label{sec:mechanism:layernorm}

A LayerNorm rescales each token to unit variance, so additive noise that raises the variance is partly divided away and a mask is not. The attention input's LayerNorm passes \AbsorptionAttentionInputGaussianSurvival{} of a Gaussian disturbance and \AbsorptionAttentionInputTokenMaskSurvival{} of a token mask on \AbsorptionModels{} models, against \AbsorptionControlSurvival{} where no LayerNorm follows (\cref{app:layernorm}). The adaptive rule therefore picks a larger rate for noise, and noise injected after the MLP's LayerNorm escapes the absorption and reads \StaircaseOperatorsBeforeMlpGaussianAuroc{}, the apparent reversal of \cref{sec:results:operator}.

\subsection{The original account holds for some attacks}
\label{sec:mechanism:phenomenon}

The PSBD paper explains its statistic by neuron bias, under which clean inputs whose prediction shifts should land on the backdoor's target. On ViT \ShiftToTargetDriftAttacks{} drift toward the target on some dataset, while \ShiftToTargetNoDriftAttacks{} stay below \ShiftToTargetNoDriftFloor{} on every dataset (\cref{app:shift-target}). PSBD-TM detects Blend and TaCT, which show no drift, so neuron bias describes some attacks rather than the statistic's mechanism.
